\documentclass[12pt, a4paper]{article}

% Языки и шрифты (XeLaTeX)
\usepackage{polyglossia}
\setdefaultlanguage{russian}
\setotherlanguage{english}
\setmainfont{Times New Roman}[Ligatures=TeX]
\setsansfont{Arial}
\setmonofont{Courier New}

% Улучшение типографики
\usepackage{microtype}

% Математика
\usepackage{amsmath, amssymb, amsthm}
\usepackage{mathtools}
\usepackage{bm}

% Оформление теорем
\theoremstyle{plain}
\newtheorem{theorem}{Теорема}[section]
\newtheorem{proposition}[theorem]{Утверждение}
\newtheorem{lemma}[theorem]{Лемма}
\newtheorem{corollary}[theorem]{Следствие}

\theoremstyle{definition}
\newtheorem{definition}[theorem]{Определение}
\newtheorem{example}[theorem]{Пример}

\theoremstyle{remark}
\newtheorem{remark}[theorem]{Замечание}

% Удобные команды и операторы
\DeclareMathOperator{\spn}{span}
\DeclareMathOperator{\diag}{diag}
\DeclareMathOperator{\rank}{rank}
\DeclareMathOperator{\Ker}{Ker}
\DeclareMathOperator{\tr}{tr}
\DeclareMathOperator{\wal}{wal} % Функции Уолша
\DeclareMathOperator{\rad}{rad} % Радикальное обращение

\newcommand{\eb}{\mathrm{e}}             
\newcommand{\ii}{\mathrm{i}}            

% Множества и числовые системы
\newcommand{\NN}{\mathbb{N}}
\newcommand{\NNz}{\mathbb{N}_0}
\newcommand{\ZZ}{\mathbb{Z}}
\newcommand{\RR}{\mathbb{R}}
\newcommand{\CC}{\mathbb{C}}
\newcommand{\FF}{\mathbb{F}}      
\newcommand{\ZZb}{\mathbb{Z}_b}    
\newcommand{\II}{[0, 1)}  
\newcommand{\PP}{\mathcal{P}} 
\newcommand{\ScrS}{\mathcal{S}}         
\newcommand{\Is}{\mathcal{I}_s}
\newcommand{\cPu}{\mathcal{P}_u}         
\newcommand{\Gbm}{G_{b,m}}             

% Векторы и матрицы
\newcommand{\va}{\boldsymbol{a}}
\newcommand{\vb}{\boldsymbol{b}}
\newcommand{\vc}{\boldsymbol{c}}
\newcommand{\vx}{\boldsymbol{x}}
\newcommand{\vy}{\boldsymbol{y}}
\newcommand{\vz}{\boldsymbol{z}}
\newcommand{\vu}{\boldsymbol{u}}
\newcommand{\vv}{\boldsymbol{v}}
\newcommand{\vw}{\boldsymbol{w}}
\newcommand{\vg}{\boldsymbol{g}}   % Столбцы генерирующей матрицы
\newcommand{\vh}{\boldsymbol{h}}   % Вектор частот (для рядов Фурье/Уолша)
\newcommand{\vzero}{\boldsymbol{0}}
\newcommand{\vone}{\boldsymbol{1}}
\newcommand{\ve}{\boldsymbol{e}}  

% Прочие обозначения
\newcommand{\ip}[2]{\langle #1,\, #2 \rangle}
\newcommand{\abs}[1]{\left| #1 \right|}
\newcommand{\norm}[1]{\left\| #1 \right\|}
\newcommand{\floor}[1]{\lfloor #1 \rfloor}
\newcommand{\ceil}[1]{\lceil #1 \rceil}

% Геометрия страницы
\usepackage[left=25mm, right=25mm, top=25mm, bottom=25mm]{geometry}

% Графика
\usepackage{graphicx}

% Гиперссылки
\usepackage{hyperref}
\hypersetup{
    colorlinks=true, 
    linkcolor=black,      
    citecolor=blue,       
    urlcolor=blue,
    pdftitle={Введение в теорию цифровых сетей и последовательностей},
    pdfauthor={В. А. Белов}
}

% Библиография
\usepackage[backend=biber, style=numeric, sorting=nyt]{biblatex}
\addbibresource{../../literature/bibliography.bib}

\title{Введение в теорию цифровых сетей и последовательностей}
\author{В. А. Белов \\ \small Российский технологический университет МИРЭА}
\date{\today}

\begin{document}

\maketitle

\begin{abstract}
Настоящий документ служит введением в теорию цифровых сетей и последовательностей, подробно описанных в фундаментальной монографии Йозефа Дика и Фридриха Пиллихшаммера \cite{dick2010digital}.
\end{abstract}

\section{Обозначения}

В работе используется единая система обозначений, представленная ниже:

\subsection*{Множества и числовые системы}
\begin{description}
    \item[$\NN$] - Множество натуральных чисел.
    \item[$\NNz$] - Множество натуральных чисел, включая ноль.
    \item[$\ZZ$] - Кольцо целых чисел.
    \item[$\RR$] - Поле действительных чисел.
    \item[$\CC$] - Поле комплексных чисел.
    \item[$\ZZb$] - Кольцо вычетов по модулю $b$ (отождествляется с $\{0, \dots, b-1\}$ с операциями сложения и умножения по модулю $b$).
    \item[$\FF_b$] - Конечное поле из $b$ элементов ($b=p^k$, $p$ - простое). При $b=p$ отождествляется с $\ZZb$. Элементы иногда обозначаются в виде $\overline{0}, \overline{1}, \dots, \overline{b-1}$.
    \item[$|X|$] - Мощность множества $X$.
    \item[$X^m$] - $m$-кратное декартово произведение множества $X$.
    \item[$(X^m)^\top$] - Множество $m$-мерных столбцовых векторов над $X$.
    \item[$\PP$] - Конечное точечное множество в $\II^s$. Интерпретируется как «мультимножество» (набор элементов, где важна кратность появления каждой точки)..
    \item[$\ScrS$] - Бесконечная последовательность точек в $\II^s$.
    \item[$\Is$] - Индексное множество $\{1, \dots, s\}$.
    \item[$u, v, \dots$] - Подмножества индексного множества $\Is$.
    \item[$\cPu$] - Точечное множество в $[0,1]^{|u|}$, полученное проекцией точек из $\PP$ на координаты, соответствующие $u \subseteq \Is$.
    \item[{$\FF_b((x^{-1})), \ZZb((x^{-1}))$}] - Поле формальных рядов Лорана над $\FF_b$ или $\ZZb$.
    \item[$G_{b,m}$] = $\{q \in \FF_b[x] : \deg(q) < m\}$.
    \item[$\gamma$] - Набор неотрицательных весов $\gamma = \{\gamma_u : u \subseteq \Is\}$. В случае весов произведения $\gamma = (\gamma_i)_{i \ge 1}$ и $\gamma_u = \prod_{i \in u} \gamma_i$.
\end{description}

\subsection*{Векторы и матрицы}
Векторы рассматриваются как столбцы. Жирное начертание обозначает векторы.
\begin{description}
    \item[$\vx, \vy, \vz$] - Векторы в $\RR^s$ или $\FF_b^s$.
    \item[$\vx_u$] - Проекция вектора $\vx$ на координаты с индексами из $u \subseteq \Is$.
    \item[$( \vx_u, \vone )$] - Вектор, у которого компоненты с индексами из $u$ взяты от $\vx$, а остальные равны $1$.
    \item[$( \vx_u, \vzero )$] - Аналогично, но остальные компоненты равны $0$.
    \item[$( \vx_u, \vw )$] - Вектор, у которого $i$-я компонента равна $x_i$, если $i \in u$, и $w_i$ в противном случае.
    \item[$A, B, C, \dots$] - Матрицы над $\FF_b$ размера $m \times m$ или $\NN \times \NN$.
    \item[$C^\top$] - Транспонированная матрица.
    \item[$C^{(m)}$] - Левый верхний $m \times m$-блок матрицы $C$.
    \item[$C^{(m \times n)}$] - Левый верхний $m \times n$-блок матрицы $C$.
\end{description}

\subsection*{Алгебраические структуры и функции}
\begin{description}
    \item[$\mathrm{i}$] - Мнимая единица ($\mathrm{i}^2 = -1$).
    \item[$\omega_b$] = $\eb^{2\pi \mathrm{i}/b}$.
    \item[$\mathrm{wal}_{b,k}$] - $k$-я $b$-адическая функция Уолша.
    \item[$\chi_J(\vx)$] - Характеристическая функция множества $J$: равна $1$, если $\vx \in J$, и $0$ в противном случае.
    \item[$\phi_b$] - $b$-адическая функция радикального обращения.
    \item[$\varphi$] - Биекция $\{0, \dots, b-1\} \to \FF_b$.
    \item[$\varphi^{-1}$] - Обратная биекция.
\end{description}

\subsection*{Цифровые разложения и операции}
\begin{description}
    \item[$\overrightarrow{\mathrm{tr}}_m(k)$] - Вектор первых $m$ цифр $b$-адического разложения $k$.
    \item[$\mathrm{tr}_m(k)$] - Целое число, соответствующее первым $m$ цифрам $b$-адического разложения $k$.
    \item[$\pi_m(\vc)$] - Проекция вектора $\vc \in \FF_b^N$ на первые $m$ компонент.
    \item[$d \mid n$] - $d$ делит $n$.
    \item[$\{x\}$] - Дробная часть числа $x$.
    \item[$\lfloor x \rfloor$] - Целая часть числа $x$.
    \item[$\lceil x \rceil$] - Наименьшее целое, не меньшее $x$.
    \item[$\overline{a}$] - Комплексно сопряжённое к $a$.
    \item[$|\vx|_1$] - $L_1$-норма: $|\vx|_1 = |x_1| + \dots + |x_s|$.
    \item[$|\vx|_\infty$] - Максимальная норма: $|\vx|_\infty = \max_{1\le i\le s} |x_i|$.
\end{description}

\subsection*{Параметры качества и метрики}
\begin{description}
    \item[$A(J, N, \ScrS)$] - Число точек последовательности $\ScrS$ среди первых $N$, попавших в $J$.
    \item[$A(J, N, \PP)$] - Число точек множества $\PP$, попавших в $J$.
    \item[$\lambda_s$] - $s$-мерная мера Лебега.
    \item[$D_N^*$] - Звёздная неравномерность (star discrepancy).
    \item[$D_{N,\gamma}^*$] - Взвешенная звёздная неравномерность.
    \item[$D_N$] - Экстремальная неравномерность (extreme discrepancy).
    \item[$L_{q,N}$] - $L_q$-неравномерность.
    \item[$L_{q,N,\gamma}$] - Взвешенная $L_q$-неравномерность.
    \item[$B_k$] - $k$-й многочлен Бернулли.
    \item[$I(f)$] - Интеграл функции $f$ по $\II^s$.
    \item[$Q_N(f)$] - Правило квази-Монте-Карло: $Q_N(f) = \frac{1}{N} \sum_{n=0}^{N-1} f(\vx_n)$.
    \item[$\mathrm{Prob}$] - Вероятность.
    \item[$\mathbb{E}$] - Математическое ожидание.
    \item[$\mathrm{Var}$] - Дисперсия.
    \item[$O(f(N))$] - Большое $O$-обозначение.
\end{description}

\printbibliography[heading=bibintoc, title={Список литературы}]

\end{document}